\begin{figure*}[t]
\centering
\hfil\hfil%
\begin{subfigure}[t]{.27\textwidth}
\centering
\scalebox{0.7}{%
\begin{tikzpicture}[
 node distance=2.5em, auto,
 lat/.style={draw=black, circle, minimum size=2em},
 obs/.style={circle, draw=black, fill=black!20, minimum size=2em},
 gen/.style={->, -{Stealth[length=.5em, inset=0pt]}},
 inf/.style={dashed, ->, -{Stealth[length=.5em, inset=0pt]}},
 div/.style={decorate, decoration={snake, amplitude=.15em, segment length=.8em, post length=.45em}, ->, -{Stealth[length=.5em, inset=0pt]}},
]
% Observations
\node[obs, inner sep=.02em] (o1) {$o_1,r_1$};
\node[obs, right=of o1, inner sep=.02em] (o2) {$o_2,r_2$};
\node[obs, right=of o2, inner sep=.02em] (o3) {$o_3,r_3$};
% Latents
\node[lat, above=of o1] (q11) {$s_{1|1}$};
\node[lat, above=of o2] (q22) {$s_{2|2}$};
\node[lat, above=of o3] (q33) {$s_{3|3}$};
\node[lat, above=of q22] (q21) {$s_{2|1}$};
\node[lat, above=of q33] (q32) {$s_{3|2}$};
% Generative
\path (q11) edge[gen] node {} (o1);
\path (q22) edge[gen] node {} (o2);
\path (q33) edge[gen] node {} (o3);
\path (q11) edge[gen] node {} (q21);
\path (q22) edge[gen] node {} (q32);
% Inference
\path (o1) edge[inf, bend left=40] node {} (q11);
\path (o2) edge[inf, bend left=40] node {} (q22);
\path (o3) edge[inf, bend left=40] node {} (q33);
\path (q11) edge[inf] node {} (q22);
\path (q22) edge[inf] node {} (q33);
% Divergence
\path (q22) edge[div] node {} (q21);
\path (q33) edge[div] node {} (q32);
\end{tikzpicture}}
\caption{标准变分下界}
\label{fig:elbo}
\end{subfigure}\hfil%
%
\begin{subfigure}[t]{.27\textwidth}
\centering
\scalebox{0.7}{%
\begin{tikzpicture}[
 node distance=2.5em, auto,
 lat/.style={draw=black, circle, minimum size=2em},
 obs/.style={circle, draw=black, fill=black!20, minimum size=2em},
 gen/.style={->, -{Stealth[length=.5em, inset=0pt]}},
 inf/.style={dashed, ->, -{Stealth[length=.5em, inset=0pt]}},
 div/.style={decorate, decoration={snake, amplitude=.15em, segment length=.8em, post length=.45em}, ->, -{Stealth[length=.5em, inset=0pt]}},
]
% Observations
\node[obs, inner sep=.02em] (o1) {$o_1,r_1$};
\node[obs, right=of o1, inner sep=.02em] (o2) {$o_2,r_2$};
\node[obs, right=of o2, inner sep=.02em] (o3) {$o_3,r_3$};
% Latents
\node[lat, above=of o1] (q11) {$s_{1|1}$};
\node[lat, above=of o2] (q22) {$s_{2|2}$};
\node[lat, above=of o3] (q33) {$s_{3|3}$};
\node[lat, above=of q22] (q21) {$s_{2|1}$};
\node[lat, above=of q33] (q32) {$s_{3|2}$};
\node[lat, above=of q32] (q31) {$s_{3|1}$};
% Generative
\path (q11) edge[gen] node {} (o1);
\path (q22) edge[gen] node {} (o2);
\path (q33) edge[gen] node {} (o3);
\path (q11) edge[gen] node {} (q21);
\path (q22) edge[gen] node {} (q32);
\path (q21) edge[gen] node {} (q31);
\path (q21) edge[gen, bend left=30] node {} (o2);
\path (q32) edge[gen, bend left=30] node {} (o3);
\path (q31) edge[gen, bend left=50] node {} (o3);
% Inference
\path (o1) edge[inf, bend left=40] node {} (q11);
\path (o2) edge[inf, bend left=40] node {} (q22);
\path (o3) edge[inf, bend left=40] node {} (q33);
\path (q11) edge[inf] node {} (q22);
\path (q22) edge[inf] node {} (q33);
\end{tikzpicture}}
\caption{观测过冲}
\label{fig:obsov}
\end{subfigure}\hfil%
%
\begin{subfigure}[t]{.27\textwidth}
\centering
\scalebox{0.7}{%
\begin{tikzpicture}[
 node distance=2.5em, auto,
 lat/.style={draw=black, circle, minimum size=2em},
 obs/.style={circle, draw=black, fill=black!20, minimum size=2em},
 gen/.style={->, -{Stealth[length=.5em, inset=0pt]}},
 inf/.style={dashed, ->, -{Stealth[length=.5em, inset=0pt]}},
 div/.style={decorate, decoration={snake, amplitude=.15em, segment length=.8em, post length=.45em}, ->, -{Stealth[length=.5em, inset=0pt]}},
]
% Observations
\node[obs, inner sep=.02em] (o1) {$o_1,r_1$};
\node[obs, right=of o1, inner sep=.02em] (o2) {$o_2,r_2$};
\node[obs, right=of o2, inner sep=.02em] (o3) {$o_3,r_3$};
% Latents
\node[lat, above=of o1] (q11) {$s_{1|1}$};
\node[lat, above=of o2] (q22) {$s_{2|2}$};
\node[lat, above=of o3] (q33) {$s_{3|3}$};
\node[lat, above=of q22] (q21) {$s_{2|1}$};
\node[lat, above=of q33] (q32) {$s_{3|2}$};
\node[lat, above=of q32] (q31) {$s_{3|1}$};
% Generative
\path (q11) edge[gen] node {} (o1);
\path (q22) edge[gen] node {} (o2);
\path (q33) edge[gen] node {} (o3);
\path (q11) edge[gen] node {} (q21);
\path (q22) edge[gen] node {} (q32);
\path (q21) edge[gen] node {} (q31);
% Inference
\path (o1) edge[inf, bend left=40] node {} (q11);
\path (o2) edge[inf, bend left=40] node {} (q22);
\path (o3) edge[inf, bend left=40] node {} (q33);
\path (q11) edge[inf] node {} (q22);
\path (q22) edge[inf] node {} (q33);
% Divergence
\path (q22) edge[div] node {} (q21);
\path (q33) edge[div] node {} (q32);
\path (q33) edge[div, bend right=40] node {} (q31);
\end{tikzpicture}}
\caption{潜在过冲}
\label{fig:latov}
\end{subfigure}%
\hfil\hfil%
\caption{展开方案。标签 $s_{i|j}$ 表示时间步 $i$ 的状态，并以截至时间步 $j$ 的观测为条件。指向灰色圆点的箭头表示对数似然损失项，波浪箭头表示 KL 散度损失项。(a) 标准变分目标在每个时间步解码后验以计算重建损失，并在对应的先验与后验之间加入 KL 项，从而训练单步预测的转移函数。(b) 观测过冲\citep{amos2018awareness}会解码所有多步预测，以施加额外重建损失；在图像域中，这通常代价过高。(c) 潜在过冲预测所有多步先验，并在潜在空间中使这些状态信念接近相应后验，从而鼓励准确的多步预测。}
\label{fig:overshooting}
\end{figure*}
